\chapter*{Streszczenie}
\addcontentsline{toc}{chapter}{Streszczenie}

Tytuł pracy dyplomowej: Łączenie grafowych sieci neuronowych z empirycznymi 
zależnościami statystycznymi na farmakoterapeutycznym grafie wiedzy

\noindent
Przedmiotem pracy dyplomowej jest przetestowanie możliwości łączenia farmakoterapeutycznej 
wiedzy zakodowanej w formie grafu wiedzy (ang.\ knowledge graph) 
z danymi empirycznymi pacjentów. 
Zagadnienie to jest rozpatrywane na dwóch uzupełniających się zadaniach
w docelowym środowisku klinicznym.
Zadanie~A koncentruje się na predykcji krawędzi (ang.\ link prediction) w grafie mechanistycznym, 
z wykorzystaniem topologii grafu wiedzy oraz statystycznej wiedzy na temat zjawiska uzyskanej z 
danych pacjenckich.
Natomiast Zadanie~B ma wykorzystywać graf wiedzy do predykcji wystąpienia zdarzeń klinicznych, 
rozumianych jako zdarzenia wymagające interwencji klinicznej, wykorzystując graf jako 
wspólny szkielet strukturalny, podczas gdy obserwacje dotyczące poszczególnych pacjentów stanowią 
dane wejściowe na poziomie węzłów.

Oba zadania wykorzystują grafowe sieci neuronowe (ang.\ Graph Neural Networks) oraz dane empiryczne 
przekazane do modelu w formie statystyk, ze szczególnym uwzględnieniem oceny siły zależności między zmiennymi.
Implementacje w obu zadaniach różnią się docelowymi architekturami, głębokością message passingu, 
połączeniami rezydualnymi, a także sposobem wykorzystania gałęzi statystycznej w architekturze sieci.
W części statystycznej wspólnym mianownikiem jest jednak wykorzystanie metody inspirowanej Hierarchical Correlation 
Reconstruction (HCR), która umożliwia opis zależności pomiędzy zmiennymi różnych typów za pomocą 
ortonormalnych funkcji bazowych i odpowiadających im współczynników.
%, pozwalając uwzględniać zarówno proste, 
%jak i bardziej złożone struktury zależności.

Na syntetycznym grafie farmakoterapeutycznym finalny model Zadania~A wykazuje bardzo dobre wyniki 
w oficjalnym protokole pięciu scenariuszy, a ablacja pokazuje, że jawny kontekst statystyczny 
istotnie poprawia predykcję krawędzi względem samej reprezentacji grafu.
W Zadaniu~B moduł statystyczny przynosi największe korzyści w sytuacjach, gdy głębsze 
warstwy message passingu prowadzą do utraty informacji predykcyjnych, 
podczas gdy połączenia rezydualne pozwalają zachować te informacje w obrębie ścieżki grafowej.
%W rezultacie na grafie syntetycznym działanie obu mechanizmów nie zawsze się sumuje.

Ogólny charakter tych obserwacji zweryfikowano dodatkowo na rzeczywistych zbiorach międzydomenowych.
W przypadku Last-FM* (z korektą wycieku danych) stopniowe dodawanie kontekstu 
specyficznego dla kandydata systematycznie poprawia jakość rankingu względem samego enkodera grafu 
heterogenicznego.
W eksperymencie predykcji punktów końcowych na Hetionet, gałąź statystyczna 
ponownie okazuje się szczególnie przydatna przy głębszych warsztwach modeli grafowych.
Na heterogenicznym grafie transkacji kartowych z znacznikiem fraud
 H2GB IEEE-CIS-G, połączenia rezydualne stabilizują uczenie, a gałąź statystyczna dodatkowo 
 poprawia detekcję oszustw względem samego message passingu z połączeniami rezydualnymi, 
 przy strukturze grafu odmiennej od syntetycznego szkieletu.

Wyniki potwierdzają wspólny wniosek metodologiczny: 
ustrukturyzowane reprezentacje grafów i jawna informacja o zależnościach 
empirycznych stanowią uzupełniający się zestaw źródeł informacji zarówno do rekonstrukcji grafów, 
jak i predykcji na poziomie pacjenta.
Ich względny wkład zależy od zadania oraz od tego, jak skutecznie grafowa sieć neuronowa zachowuje 
istotne informacje podczas kolejnych etapów propagacji i agregacji informacji w grafie.
Eksperymenty na zbiroach rzeczywistych i międzydomenowych są metodologicznymi testami obciążeniowymi, 
jednak ich wyniki należy interpretować ostrożnie przed jakimkolwiek zastosowaniem w realnych 
warunkach klinicznych.

\medskip
\noindent\textbf{Słowa kluczowe:}
grafy wiedzy,
grafowe sieci neuronowe,
grafy heterogeniczne,
predykcja linków,
predykcja klinicznych punktów końcowych,
Hierarchical Correlation Reconstruction,
farmakoterapia.
